% ============================================================
%  Póster científico minimalista (tikzposter)
%  Compilar con: pdflatex poster_minimalista.tex  (2 pasadas)
%  Tamaño: 90 x 120 cm (ancho x alto, vertical), 3 columnas
%  Para cambiar el color de acento: edita \definecolor{acento}
% ============================================================
\documentclass[25pt, a0paper, portrait,
  margin=0mm, innermargin=18mm,
  blockverticalspace=14mm, colspace=16mm, subcolspace=8mm]{tikzposter}

% ---------- Tamaño personalizado: 90 x 120 cm (ancho x alto) ----------
% tikzposter solo acepta a0/a1/a2 como opción; aquí se sobrescribe con geometry.
\geometry{paperwidth=90cm, paperheight=120cm, margin=0mm}

\usepackage[utf8]{inputenc}
\usepackage[T1]{fontenc}
\usepackage[spanish]{babel}
\usepackage{lmodern}
\renewcommand{\familydefault}{\sfdefault}
\usepackage{amsmath, amssymb}
\usepackage{graphicx}
\usepackage{booktabs}
\usepackage{enumitem}
\usepackage{hyperref}

% ---------- Paleta (minimalista: negro, gris y un acento) ----------
\definecolor{tinta}{HTML}{1A1A1A}
\definecolor{gris}{HTML}{6B6B6B}
\definecolor{acento}{HTML}{0057B8}

\definecolorstyle{Minimal}{}{
  \colorlet{backgroundcolor}{white}
  \colorlet{framecolor}{white}
  \colorlet{titlefgcolor}{tinta}
  \colorlet{titlebgcolor}{white}
  \colorlet{blocktitlebgcolor}{white}
  \colorlet{blocktitlefgcolor}{acento}
  \colorlet{blockbodybgcolor}{white}
  \colorlet{blockbodyfgcolor}{tinta}
  \colorlet{innerblocktitlebgcolor}{white}
  \colorlet{innerblocktitlefgcolor}{acento}
  \colorlet{innerblockbodybgcolor}{white}
  \colorlet{innerblockbodyfgcolor}{tinta}
  \colorlet{notefgcolor}{tinta}
  \colorlet{notebgcolor}{white}
  \colorlet{noteframecolor}{acento}
}

% ---------- Título: sin caja ni sombras ----------
\definetitlestyle{Minimal}{
  width=\textwidth, roundedcorners=0, linewidth=0pt, innersep=0pt,
  titletotopverticalspace=18mm, titletoblockverticalspace=14mm
}{}

% ---------- Bloques: solo una línea fina bajo el encabezado ----------
\defineblockstyle{Minimal}{
  titlewidthscale=1, bodywidthscale=1, titleleft,
  titleoffsetx=0pt, titleoffsety=0pt,
  bodyoffsetx=0pt, bodyoffsety=0pt,
  bodyverticalshift=8mm, roundedcorners=0, linewidth=0pt,
  titleinnersep=4mm, bodyinnersep=0pt
}{
  \draw[color=blocktitlefgcolor, line width=1.5pt]
    (blocktitle.south west) -- (blocktitle.south east);
}

\definebackgroundstyle{Minimal}{}
\usecolorstyle{Minimal}
\usetitlestyle{Minimal}
\useblockstyle{Minimal}
\usebackgroundstyle{Minimal}

% Estilo del título principal (edita el texto del título en las 3 líneas de abajo)
\makeatletter
\settitle{\centering\parbox{80cm}{\centering\color{titlefgcolor}%
  {\bfseries\fontsize{88}{100}\selectfont
    Desarrollo asistido por Inteligencia Artificial\\
    de un resolvedor de sistemas de ecuaciones\\
    complejas para análisis de circuitos \par}\vspace{12mm}
  {\fontsize{52}{62}\selectfont \@author \par}\vspace{6mm}
  {\color{gris}\fontsize{42}{52}\selectfont \@institute \par}}}
\makeatother

% Tamaños de texto (póster 90x120 cm). Ajusta el cuerpo con \large.
\makeatletter
\renewcommand\footnotesize{\@setfontsize\footnotesize{32}{40}}
\renewcommand\small{\@setfontsize\small{38}{48}}
\renewcommand\normalsize{\@setfontsize\normalsize{42}{54}}
\renewcommand\large{\@setfontsize\large{48}{60}}      % cuerpo de los bloques
\renewcommand\LARGE{\@setfontsize\LARGE{64}{76}}      % títulos de bloque
\makeatother

% Listas compactas
\setlist{itemsep=2mm, topsep=2mm, leftmargin=*}

% ---------- Carpetas donde buscar las figuras ----------
\graphicspath{{./}{Complex calc/}{Projects/Complex Calc/LaTeXProject/Complex calc/}}

% \figuraopt{archivo}{fracción de \linewidth}{pie de figura}
% Si el archivo no existe, dibuja un recuadro gris en su lugar.
\newcommand{\figuraopt}[3]{%
  \begin{tikzfigure}[#3]
    \IfFileExists{#1}{\includegraphics[width=#2\linewidth]{#1}}{%
      \tikz\draw[gris, line width=1pt] (0,0) rectangle (#2\linewidth, 0.5*#2\linewidth)
        node[midway, gris]{\fontsize{30}{36}\selectfont #1};}%
  \end{tikzfigure}}

% ---------- Metadatos ----------
% (El texto del título que se muestra se edita dentro de \settitle, arriba, con saltos \\ manuales.)
\title{Desarrollo asistido por IA de un resolvedor de sistemas de ecuaciones complejas}
\author{Orlando Contreras Reyes, Iker Emiliano García Barrera,\\
        Dr. Luís Alejandro Flores Oropeza, Dr. Alejandro Román Loera}
\institute{Universidad Autónoma de Aguascalientes \quad
           \texttt{al348390@edu.uaa.mx}}

\begin{document}
\tikzposterlatexaffectionproofoff  % quita el logo "LaTeX TikZposter" del pie
\maketitle

\begin{columns}

% ======================= COLUMNA 1 =======================
\column{0.333}

\block{Introducción}{
  El análisis de circuitos (mallas o nodos) exige resolver sistemas de
  ecuaciones con números complejos. Las soluciones habituales presentan
  limitaciones:

  \begin{itemize}
    \item \textbf{Regla de Cramer:} método poco eficiente.
    \item \textbf{Calculadoras físicas y simuladores} (p.\,ej. HP-48G, Droid~48):
          ingreso de datos poco eficiente y resultados solo en forma
          rectangular o solo polar.
    \item \textbf{Deshonestidad académica:} uso de herramientas sin
          comprender el proceso.
  \end{itemize}

  \vspace{4mm}
  \textbf{Objetivo.} Crear una aplicación de escritorio, portátil y
  gratuita que resuelva sistemas de ecuaciones complejas y muestre los
  resultados en ambas formas, desarrollada junto con IA.
}

\block{Métodos: herramientas}{
  \begin{itemize}
    \item \textbf{Lenguaje y entorno:} Python 3.10+, Visual Studio Code.
    \item \texttt{NumPy}: álgebra lineal compleja.
    \item \texttt{CustomTkinter} y \texttt{Tkinter}: interfaz gráfica.
    \item \texttt{Pillow}, \texttt{OS} y \texttt{JSON}: recursos gráficos,
          archivos e historial de matrices.
    \item \textbf{IA:} ChatGPT (OpenAI) para código base, refactorización
          y diseño de la GUI.
    \item \texttt{PyInstaller}: ejecutable único
          (\texttt{--onefile}, \texttt{--noconsole}) para Windows y GNU/Linux.
  \end{itemize}
}

\block{Filosofía}{
  \textit{Open Source} y \textit{Vibe Coding}: el proyecto lo realizaron
  alumnos con bases de programación en C para microcontroladores. Con
  buenas bases, el ingeniero puede construir sus propias herramientas.
  Se publica bajo licencia MIT para que otros estudiantes aprendan,
  contribuyan y la reutilicen.
}

% ======================= COLUMNA 2 =======================
\column{0.333}

\block{Métodos: arquitectura y algoritmo}{
  \figuraopt{src/Architecture.png}{0.95}{Arquitectura modular.}

  \vspace{4mm}
  Un sistema de $n \times n$ se plantea como
  \begin{equation*}
    [\mathbf{Z}][\mathbf{I}] = [\mathbf{V}]
    \;\Rightarrow\;
    [\mathbf{I}] = [\mathbf{Z}]^{-1}[\mathbf{V}]
  \end{equation*}
  y se resuelve con \texttt{numpy.linalg.solve}. Los resultados se
  convierten a forma polar:
  \begin{equation*}
    r = \sqrt{a^2 + b^2}, \qquad \theta = \arctan\!\left(\frac{b}{a}\right)
  \end{equation*}
  Entrada en forma rectangular ($a + jb$) o polar ($r\angle\theta^\circ$).
}

\block{Flujo de trabajo con IA}{
  \begin{enumerate}
    \item Maquetas iniciales de ventanas y celdas matriciales editables.
    \item Solución de errores con arreglos complejos en tiempo real.
    \item Refactorización: separar cálculo y renderizado.
    \item \textbf{Validación humana:} cada fragmento fue revisado,
          probado y contrastado con ejercicios resueltos a mano.
  \end{enumerate}
}

\block{Discusión}{
  \begin{itemize}
    \item Útil en Análisis de Circuitos, Electrónica, Señales y Sistemas
          y Métodos Matemáticos.
    \item Historial de operaciones y guardado/carga de matrices.
    \item Interfaz reorganizada para ser clara a usuarios sin experiencia.
    \item Tkinter limita los temas visuales; se personalizó un archivo
          \texttt{.json} de estilos.
  \end{itemize}
}

% ======================= COLUMNA 3 =======================
\column{0.333}

\block{Resultados}{
  \figuraopt{RectSol.png}{0.95}{Solución de un sistema complejo en forma rectangular.}

  \vspace{4mm}
  \figuraopt{SolDC.png}{0.95}{Solución de un sistema complejo en forma polar.}
}

\block{Conclusiones}{
  \begin{enumerate}
    \item La IA es un apoyo, no un sustituto del conocimiento del
          programador: reduce el tiempo de desarrollo, pero exige
          entender y validar el código.
    \item Con bases de C para microcontroladores se logró una
          aplicación funcional con entrada y conversión rectangular/polar.
    \item El ejecutable facilita su distribución sin configurar el entorno.
    \item Licencia MIT: base abierta para mejoras y herramientas similares.
  \end{enumerate}
}

\block{Referencias}{
  \footnotesize
  \begin{enumerate}[label={[\arabic*]}, itemsep=1mm]
    \item H. Anton y C. Rorres, \textit{Álgebra lineal con aplicaciones}. Wiley, 2013.
    \item D. C. Lay, \textit{Álgebra lineal y sus aplicaciones}, 4.ª ed. Pearson, 2012.
    \item NumPy Developers, \textit{numpy.linalg.solve}. NumPy Documentation, 2026.
    \item Python Software Foundation, \textit{Built-in Types} y \textit{tkinter}. Python Documentation, 2026.
    \item D. Cortesi \textit{et al.}, \textit{PyInstaller Manual}, 2026.
    \item OpenAI, \textit{ChatGPT}, 2026.
  \end{enumerate}
}

\block{Contacto y código}{
  \footnotesize
  \texttt{al348390@edu.uaa.mx} \\
  \url{https://github.com/DasReyxr/Py-ComplexCalc}
  % \vspace{4mm}
  % \includegraphics[width=0.25\linewidth]{qr.png}
}

\end{columns}

\end{document}